\section{Introduction}
\label{sec:introduction}

The TPC-18 stopping argument isolates a primitive
large-determinant/low-gcd residual inside each admissible symmetric
tail block \citep{WangTPC18}.  Its conditional global conclusion
requires more than a good estimate on each block: the physical hard
packet must be represented by the entire proposed block family, with
the correct unit coefficients, overlaps, masks, and normalization, up
to a controlled \(o(X)\) physical remainder.  That requirement was
deliberately left as an all-block representation hypothesis.

A bounded-overlap partition is often enough for inequalities because
one may estimate each atom a bounded number of times.  It is not
enough for an identity.  If two blocks overlap, the shared atom
receives the sum of both block coefficients.  If the block weights
cannot reproduce the target coefficient exactly, dividing by a
support multiplicity does not repair signs, phases, masks, or
nonconstant block profiles.  This distinction is especially important
when a later argument is calibrated to a single global main term.
An exact coefficient identity is a sufficient, zero-remainder way to
verify this hypothesis.  It is stronger than the original asymptotic
statement: a nonzero coefficient residual may also be admissible if a
separate theorem proves that its physical image is \(o(X)\).

The purpose of this paper is to turn that qualitative warning into an
exact finite-dimensional certificate.  Let \(\Omega_X\) be the actual
atom set and let \(w_X\) be the literal coefficient of the physical
packet.  Each candidate block contributes one exact column \(b_j\).
The synthesis matrix
\[
 B_X=(b_1\ \cdots\ b_m)
\]
contains all overlap information.  A zero-remainder cover is precisely
\[
 B_Xc_X=w_X,
\]
together with the coefficient restrictions demanded by the
decomposition.  Linear algebra then separates three questions:
\begin{enumerate}[label=(\roman*),leftmargin=2.2em]
\item does the target lie in the column space?
\item if signed coefficients are forbidden, does it lie in the
  allowed cone or polytope?
\item if it does, what coefficient amplification is paid when the
  block estimates are reassembled?
\end{enumerate}

\subsection{Main results}

The first result is the exact subspace certificate.  With
\(P_B=BB^\dagger\),
\[
 \inf_c\|w-Bc\|_2=\|(I-P_B)w\|_2.
\]
The residual is zero exactly when a cover exists.  When it is nonzero,
its normalized direction annihilates every candidate block while
detecting the target, so it is a short dual certificate of failure.

The second result is the constrained analogue.  For real
nonnegative coefficients, a cover exists exactly when
\[
 B^{\mathsf T}y\ge0\quad\Longrightarrow\quad
 \langle y,w\rangle\ge0.
\]
Failure therefore has a Farkas witness.  Cone projection yields the
quantitative distance formula
\[
 \dist(w,B\R_+^m)
 =
 \max_{\substack{\|y\|_2\le1\\B^{\mathsf T}y\le0}}
 \langle y,w\rangle.
\]
A support-function formula handles any closed convex allowed
coefficient image.

The third result consists of sharp countermodels.  A disjoint
partition of \(2n\) atoms into \(n\) two-point blocks has overlap one,
covers every atom as a set, and yet may have relative least-squares
defect exactly \(1/\sqrt{10}\) for every \(n\).  A separate two-column
example has an exact signed cover but no nonnegative cover, with cone
distance \(1/\sqrt2\).  Thus neither bounded overlap nor full set
coverage implies exact or asymptotically approximate coefficient
coverage.

Finally, we state the only automatic normalization transfer.  If
\(Bc=w\) exactly, then every linear physical statistic with its
original normalizer satisfies the exact reassembly identity, and the
block error is amplified by \(\|c\|_1\).  If
\(w=Bc+r\) only approximately, the physical remainder is the
additional term carried by \(r\).  Small Euclidean residual alone
does not remove it, although a separate physical bound on that term
can verify the original asymptotic TPC-18 route.

\subsection{Claim levels}

The projection identities, dual formulas, counterexamples, and
conditioning bounds are unconditional \(\Lzero\) statements.  Their
literal placement on the TPC-18 atom space gives an exact \(\Lone\)
certificate schema.  We do not construct the physical incidence
matrix from the whole preceding normal form, prove its feasibility,
or prove the primitive determinant-dispersion estimate.  Consequently
there is no new \(\Ltwo\) arithmetic theorem.

\subsection{Organization}

\Cref{sec:literal-model} defines the literal atom and block spaces.
\Cref{sec:dual-certificates} gives exact subspace, cone, and convex
certificates.  \Cref{sec:counterexamples} proves the overlap and
coarsening obstructions.  \Cref{sec:stability-normalization} records
conditioning and normalization costs.
\Cref{sec:tpc18-crosswalk} states the precise conditional crosswalk to
TPC-18.  \Cref{sec:separate-ledgers,sec:claim-boundary} keep the
remaining gates and claim boundary explicit.
